% Options for packages loaded elsewhere
\PassOptionsToPackage{unicode}{hyperref}
\PassOptionsToPackage{hyphens}{url}
\PassOptionsToPackage{dvipsnames,svgnames,x11names}{xcolor}
%
\documentclass[
  11pt,
]{article}
\usepackage{amsmath,amssymb}
\usepackage{iftex}
\ifPDFTeX
  \usepackage[T1]{fontenc}
  \usepackage[utf8]{inputenc}
  \usepackage{textcomp} % provide euro and other symbols
\else % if luatex or xetex
  \usepackage{unicode-math} % this also loads fontspec
  \defaultfontfeatures{Scale=MatchLowercase}
  \defaultfontfeatures[\rmfamily]{Ligatures=TeX,Scale=1}
\fi
\usepackage{lmodern}
\ifPDFTeX\else
  % xetex/luatex font selection
\fi
% Use upquote if available, for straight quotes in verbatim environments
\IfFileExists{upquote.sty}{\usepackage{upquote}}{}
\IfFileExists{microtype.sty}{% use microtype if available
  \usepackage[]{microtype}
  \UseMicrotypeSet[protrusion]{basicmath} % disable protrusion for tt fonts
}{}
\makeatletter
\@ifundefined{KOMAClassName}{% if non-KOMA class
  \IfFileExists{parskip.sty}{%
    \usepackage{parskip}
  }{% else
    \setlength{\parindent}{0pt}
    \setlength{\parskip}{6pt plus 2pt minus 1pt}}
}{% if KOMA class
  \KOMAoptions{parskip=half}}
\makeatother
\usepackage{xcolor}
\usepackage[margin=1in]{geometry}
\usepackage{longtable,booktabs,array}
\usepackage{calc} % for calculating minipage widths
% Correct order of tables after \paragraph or \subparagraph
\usepackage{etoolbox}
\makeatletter
\patchcmd\longtable{\par}{\if@noskipsec\mbox{}\fi\par}{}{}
\makeatother
% Allow footnotes in longtable head/foot
\IfFileExists{footnotehyper.sty}{\usepackage{footnotehyper}}{\usepackage{footnote}}
\makesavenoteenv{longtable}
\setlength{\emergencystretch}{3em} % prevent overfull lines
\providecommand{\tightlist}{%
  \setlength{\itemsep}{0pt}\setlength{\parskip}{0pt}}
\setcounter{secnumdepth}{-\maxdimen} % remove section numbering
\ifLuaTeX
  \usepackage{selnolig}  % disable illegal ligatures
\fi
\IfFileExists{bookmark.sty}{\usepackage{bookmark}}{\usepackage{hyperref}}
\IfFileExists{xurl.sty}{\usepackage{xurl}}{} % add URL line breaks if available
\urlstyle{same}
\hypersetup{
  pdftitle={Proper-time theory applied to IBM-quantum coherence: a microwave-coupling angle/shape research arc},
  pdfauthor={Trey Morris with Claude Opus 4.7},
  pdfsubject={White paper proposing a three-phase research arc applying
Gill's dual / proper-time framework to transmon-qubit coherence on IBM
Quantum hardware, with a falsifiable Phase-B prediction tied to
microwave-coupling angle/shape modulation of the proper-time Maxwell
dissipative term},
  colorlinks=true,
  linkcolor={blue},
  filecolor={Maroon},
  citecolor={Blue},
  urlcolor={blue},
  pdfcreator={LaTeX via pandoc}}

\title{Proper-time theory applied to IBM-quantum coherence: a
microwave-coupling angle/shape research arc}
\author{Trey Morris with Claude Opus 4.7}
\date{2026-05-31}

\begin{document}
\maketitle

\hypertarget{proper-time-theory-on-ibm-quantum-coherence-microwave-coupling-angleshape-focus}{%
\section{Proper-time theory on IBM quantum coherence ---
microwave-coupling angle/shape
focus}\label{proper-time-theory-on-ibm-quantum-coherence-microwave-coupling-angleshape-focus}}

\textbf{Date:} 2026-05-31. \textbf{From:} Trey Morris (with Claude Opus
4.7). \textbf{Re:} Three-phase research arc proposal; audience default
\emph{internal-campaign planning} (see Asks). Tracks against
\href{https://github.com/temoTxt/PyPhysics/issues/103}{issue \#103}.
Builds on the dual Maxwell verification (Eq. (4), see
\href{../Equation_Verification/Two_Mathematically_Equivalent_Versions_of_Maxwells_Equations.md\#L188}{\texttt{Two\_Mathematically\_Equivalent\_Versions\_of\_Maxwells\_Equations.md}})
and the precision-atomic campaigns
\href{https://github.com/temoTxt/PyPhysics/issues/50}{\#50},
\href{https://github.com/temoTxt/PyPhysics/issues/78}{\#78},
\href{https://github.com/temoTxt/PyPhysics/issues/82}{\#82}.

\begin{center}\rule{0.5\linewidth}{0.5pt}\end{center}

\hypertarget{page-1-the-apparatus-side-hook}{%
\subsection{Page 1 --- the apparatus-side
hook}\label{page-1-the-apparatus-side-hook}}

\textbf{What is known apparatus-side.} Gill's proper-time formulation
produces a dual Maxwell wave equation, verified as Eq. (4) of
\href{../Equation_Verification/Two_Mathematically_Equivalent_Versions_of_Maxwells_Equations.md\#L188}{\texttt{Two\_Mathematically\_Equivalent\_Versions\_of\_Maxwells\_Equations.md}},
that differs from the textbook form by a single first-derivative term.
The verification doc records the discriminating statement at
\href{../Equation_Verification/Two_Mathematically_Equivalent_Versions_of_Maxwells_Equations.md\#L232}{line
232}: Eq. (4) reproduces the standard wave equation under \(c\to b\),
\(t\to\tau\), \(\mathbf{w}\to\mathbf{u}\), \textbf{plus} an extra
dissipative term whose coefficient is
\(-(\mathbf{u}\!\cdot\!\mathbf{a})/b^4\) in front of
\(\partial_\tau\mathbf{E}\) (or \(\partial_\tau\mathbf{B}\)). The same
structure resurfaces as the third term of the modified
Liénard\(\,\)--\(\,\)Wiechert fields at
\href{../Equation_Verification/Two_Mathematically_Equivalent_Versions_of_Maxwells_Equations.md\#L332}{line
332}: \emph{``the third term is new and tied to the dissipative
coefficient \((\mathbf{u}\!\cdot\!\mathbf{a})/b^4\).''} The term
vanishes identically when \(\mathbf{a}=\mathbf{0}\) (source
unaccelerated) and inherits an angular dependence through
\(\mathbf{u}\!\cdot\!\mathbf{a}=|\mathbf{u}||\mathbf{a}|\cos\theta\) and
a temporal/shape dependence through \(\mathbf{a}(\tau)\). This is the
\emph{only} structural deviation Gill's classical-field theory makes
from textbook Maxwell at this order; every other dual modification is a
\(c\to b\) substitution or a kinematic re-parametrisation.

\textbf{Why it might bear on transmon-qubit coherence.} A transmon-qubit
gate operation drives a Cooper-pair condensate with a shaped microwave
pulse \(A(\tau)\) along a specific drive axis \(\hat{\mathbf{d}}\), so
the source is \emph{accelerated by construction} during gate operations
and the framework's extra term \emph{can} contribute. The substantive
open question --- the core of Phase A --- is which condensate kinematic
quantity plays the role of \(\mathbf{u}\)/\(\mathbf{a}\) in a transmon:
candidate identifications include (i) the supercurrent velocity
\(\mathbf{u}_s\propto\nabla\varphi\) from the phase gradient, (ii) the
AC-Josephson velocity at the drive frequency, or (iii) the
time-derivative of the Cooper-pair dipole moment. The three give
different effect-size estimates: at Cooper-pair drift speed
\(\beta^4 \lesssim 10^{-30}\) (unmeasurable), at Fermi speed
\(\beta^4 \sim 10^{-12}\) (below current \(T_2\) floor by \(\sim 6\)
orders), at AC-Josephson speed it depends on drive amplitude.
\textbf{Falsifiable hypothesis (Phase B):} shaping the microwave
coupling to minimise \((\mathbf{u}\!\cdot\!\mathbf{a})/b^4\) should
produce an angle/shape signature \emph{distinct from} the standard
Abraham\(\,\)--\(\,\)Lorentz\(\,\)--\(\,\)Dirac radiation-reaction term
(\(\propto\dot{\mathbf{a}}\)); the two are in-principle separable by
varying \(\theta\) at fixed \(\dot{\mathbf{a}}\), and any measured angle
dependence not accounted for by ALD is a falsifier in the direction of
(or against) the framework.

\begin{center}\rule{0.5\linewidth}{0.5pt}\end{center}

\hypertarget{page-2-three-phase-research-arc}{%
\subsection{Page 2 --- three-phase research
arc}\label{page-2-three-phase-research-arc}}

\textbf{Phase A \(-\) apparatus mapping.} Derive the dissipative-term
contribution to the transmon-qubit effective drive Hamiltonian (or
equivalent Lindblad operator) starting from Eq. (4), \emph{after}
committing to one of the three candidate \(\mathbf{u}\)/\(\mathbf{a}\)
identifications above (or a fourth proposed by Tepper). Acceptance: a
closed-form (or numerical-but-tractable) expression for the noise
contribution as a function of \(\theta\) and \(A(\tau)\) with explicit
\(b\to c\) and \(\beta\to 0\) limits, validated numerically against a
no-framework baseline via \texttt{qiskit-dynamics} (the \texttt{Solver}
integrates the Lindblad equation under the envelope \(A(\tau)\) entered
as a time-dependent \texttt{Signal}), and a numeric \(\Delta T_2/T_2\)
at IBM Heron's nominal drive parameters (anharmonicity \(\sim 300\) MHz,
\(t_g\sim 30\) ns). \textbf{Null-result-most-likely closure:} if
\(\Delta T_2/T_2 < 10^{-3}\) at the highest computable order, or the
contribution factorises into the standard ALD account with no residual
angle dependence separable from \(\dot{\mathbf{a}}\), the program closes
honestly at Phase A.

\textbf{Phase B \(-\) concrete predictions.} Conditional on Phase A
producing a nonzero residual, specify drive angles and pulse-shape
envelopes where the dissipative-term contribution is maximised
vs.~minimised at fixed drive amplitude. Acceptance: at least one
explicit prediction of the form \emph{``varying \(\theta\) from \(0\) to
\(\pi/2\) at fixed amplitude \(A_0\) and gate time \(t_g\) changes the
framework-predicted \(T_2\) contribution by \(\Delta T_2 = \ldots\),
with a sign,''} with the predicted effect size compared against
state-of-the-art transmon \(T_2\) (\(\sim 100\)--\(300~\mu\)s on IBM
Eagle/Heron) and against the current IBM-pulse-control angular
resolution.

\textbf{Phase C \(-\) IBM Quantum test.} Specify the IBM Quantum Network
resources: a single high-coherence transmon on Eagle- or Heron-class
hardware; DRAG/cosine pulse-shape control; angular control via the
standard IQ-modulation plane. Pre-flight in simulation via
\texttt{qiskit-aer} \texttt{NoiseModel.from\_backend()} augmented with a
custom Lindblad channel for the framework's term (equivalently,
\texttt{qiskit\_dynamics.backend.DynamicsBackend} carrying the same
operator); on-hardware execution via \texttt{qiskit-ibm-runtime}
\texttt{Sampler} under \texttt{qiskit-experiments}
\texttt{T2Hahn}/\texttt{T2Ramsey} protocols, with custom drive shapes
via \texttt{qiskit.pulse} \emph{if} pulse-level access is provisioned
(otherwise gate-level DRAG-coefficient sweeps). The discriminating
observable is the \emph{angle} dependence at fixed DD sequence, not raw
\(T_2\) --- CPMG/Uhrig suppress low-frequency dephasing and would mask
any framework contribution with overlapping spectral support, so the
framework's signature is measured as a residual after subtracting
matched-DD baselines. Acceptance verdicts: {[}OK{]} framework predicts
the measured angle/shape residual within the campaign's precision,
{[}warn{]} framework matches at precision floor only (consistent but not
discriminating), {[}X{]} framework refuted by absence of the predicted
signature at the predicted sensitivity.

\begin{center}\rule{0.5\linewidth}{0.5pt}\end{center}

\hypertarget{page-3-methods-asks-honest-scope}{%
\subsection{Page 3 --- methods, asks, honest
scope}\label{page-3-methods-asks-honest-scope}}

\textbf{Methods (incremental vs.~standard).} The table below names what
is already standard practice for transmon-coherence extension, alongside
what the framework would \emph{additionally} prescribe \emph{if} Phase A
succeeds.

\begin{longtable}[]{@{}
  >{\raggedright\arraybackslash}p{(\columnwidth - 2\tabcolsep) * \real{0.5000}}
  >{\raggedright\arraybackslash}p{(\columnwidth - 2\tabcolsep) * \real{0.5000}}@{}}
\toprule\noalign{}
\begin{minipage}[b]{\linewidth}\raggedright
Standard transmon-coherence technique
\end{minipage} & \begin{minipage}[b]{\linewidth}\raggedright
What the framework adds (if Phase A succeeds)
\end{minipage} \\
\midrule\noalign{}
\endhead
\bottomrule\noalign{}
\endlastfoot
DRAG-optimised pulse shaping (leakage suppression) &
Angle/shape-modulated drive minimising
\((\mathbf{u}\!\cdot\!\mathbf{a})/b^4\) at fixed leakage. \\
Dynamical decoupling (CPMG / Uhrig / DCG) & Proper-time-spaced CPMG; DCG
variants minimising \(\int (\mathbf{u}\!\cdot\!\mathbf{a})/b^4\,d\tau\)
at fixed fidelity. \\
Numerical baseline (\texttt{qiskit-aer} +
\texttt{NoiseModel.from\_backend()}) & Same noise model with a custom
Lindblad channel from the Phase-A operator, sim-only
\(\Delta T_2(\theta)\). \\
\end{longtable}

\textbf{Modeling + encoding pre-flight (Qiskit).} End-to-end laptop
simulation analogue before any IBM credit is burned: Phase A in
\texttt{qiskit-dynamics} (Lindblad \texttt{Solver} driven by \(A(\tau)\)
as a \texttt{Signal}); Phase B as a swept-\(\theta\) batch on
\texttt{qiskit-aer} producing the predicted \(T_2(\theta)\) ratio that
hardware will reproduce or refute; Phase C via the noise-augmented
\texttt{Aer}/\texttt{DynamicsBackend} under the same
\texttt{qiskit-experiments} protocol as the hardware run.
Encoding-sequence candidates worth flagging in parallel:
proper-time-spaced CPMG (inter-pulse interval set in \(\tau\), sensitive
to \(b/c\) at the precision floor), and DCG variants whose envelope
minimises the framework's angle/shape integral at matched gate fidelity.

\textbf{Asks (audience-dependent).} Default audience at draft-time is
\emph{internal-campaign planning}; the four candidate audiences listed
in \href{https://github.com/temoTxt/PyPhysics/issues/103}{issue \#103}
(Audience options section) are equally compatible with the body above.
Asks under the default audience: (a) Trey's sign-off on the
apparatus-mapping move on Page 1; (b) a superconducting-qubit theorist
for Phase-A vetting (identification of Cooper-pair
\(\mathbf{u}\)/\(\mathbf{a}\) in the transmon charge basis); (c)
optionally Tepper Gill or one of his students for the framework-internal
Wolfram-MCP derivation of the Phase-A expression; (d) IBM Quantum
Network compute time scoped only after Phase B produces a concrete
prediction worth burning calibration on. Cross-linked to the
framework-specs thread
\href{https://github.com/temoTxt/PyPhysics/issues/75}{\#75} and the open
proper-time one-loop frontier
\href{https://github.com/temoTxt/PyPhysics/issues/55}{\#55}.

\textbf{Honest scope (load-bearing).} This white paper proposes a
\emph{falsifiable research arc}. It does not claim a \emph{derived
prediction}. The framework's existing validation (Bethe-Salpeter
\href{https://github.com/temoTxt/PyPhysics/issues/50}{\#50}; Li
\(^{2+}\) Z-extension
\href{https://github.com/temoTxt/PyPhysics/issues/78}{\#78}; hydrogenic
g-factor Z-scan
\href{https://github.com/temoTxt/PyPhysics/issues/82}{\#82} with Outcome
C --- cutoff has no Z-dependence and the framework inherits QED's
bound-state structure) is at a precision at which the framework
reproduces textbook QED \emph{by construction}. Extending from bound
electrons in atoms to a superconducting transmon-qubit condensed-matter
device is a substantive conceptual leap, and the most likely outcome at
Phase A is \emph{null} under the explicit threshold above. In that case
the program closes honestly at Phase A and the campaign records the
closure as the framework's terminal-verdict for condensed-matter
applicability. The paper's value is in \textbf{specifying a falsifiable
test that does not currently exist in the campaign's repertoire} --- the
test is worth specifying whether or not the framework wins it, because
at present no condensed-matter prediction of the framework is on the
books at all.

--- Trey

\end{document}
